% Student paper — English, typeset text; Al2Cl6 structures as CIE clips. Compiler: XeLaTeX.
\documentclass[10pt,a4paper]{article}

\usepackage[margin=16mm,headheight=13pt,headsep=5pt,footskip=9mm]{geometry}
\usepackage{fontspec}
\usepackage[version=4]{mhchem}
\usepackage{chemfig}
\usepackage{graphicx}
\usepackage{xcolor}
\usepackage{booktabs}
\usepackage{array}
\usepackage{enumitem}
\usepackage{needspace}
\usepackage{fancyhdr}
\usepackage{amsmath,amssymb}
\usepackage{pgffor}

\raggedbottom
\setlist[enumerate]{leftmargin=1.5em,itemsep=0.12em,topsep=0.15em,parsep=0pt}
\setlist[itemize]{leftmargin=1.3em,itemsep=0.08em}

\pagestyle{fancy}
\fancyhf{}
\fancyhead[L]{\small \homeworktitle}
\fancyhead[R]{\small \homeworkmarks}
\fancyfoot[C]{\thepage}
\renewcommand{\headrulewidth}{0.3pt}

\newcommand{\homeworktitle}{Chemistry homework}
\newcommand{\homeworkmarks}{}
\newcommand{\sethomework}[2]{%
  \renewcommand{\homeworktitle}{#1}%
  \renewcommand{\homeworkmarks}{#2}%
}

\newcommand{\answerlines}[1]{%
  \par\vspace{0.08em}%
  \foreach \n in {1,...,#1}{\noindent\rule{\linewidth}{0.2pt}\\[0.62em]}%
}

\newcommand{\drawbox}[1][3.8cm]{%
  \par\vspace{0.10em}%
  \noindent\fbox{\parbox[c][#1][c]{\dimexpr\linewidth-2\fboxsep-2\fboxrule\relax}{\strut}}%
  \par
}

\graphicspath{{images/}}
\DeclareGraphicsExtensions{.png,.jpg,.jpeg,.pdf}

\setlength{\parindent}{0pt}
\setlength{\parskip}{0.10em}
\setlength{\abovedisplayskip}{0.25em}
\setlength{\belowdisplayskip}{0.25em}

\newcommand{\qhead}[2]{%
  \par\vspace{0.28em}%
  \Needspace{4.2em}%
  \noindent{\small\textbf{#1}}\hfill{\small[#2]}\par\vspace{0.08em}%
}

\newcommand{\clipq}[3]{%
  \par
  \sbox0{\includegraphics[width=0.90\linewidth,keepaspectratio]{#3}}%
  \Needspace{\dimexpr\ht0+1.25em+2pt\relax}%
  \noindent{\small\textbf{#1}}\hfill{\small[#2]}\par\vspace{0.06em}%
  \noindent\usebox0\par
}

\newcommand{\optline}[1]{%
  \par\vspace{0.06em}{\small #1}\par
}

\newcommand{\sourcefoot}[1]{%
  {\footnotesize\color{gray}#1}\par
}

\newcommand{\combokey}{%
{\small The responses \textbf{A} to \textbf{D} should be selected on the basis of}
\vspace{0.10em}
\centering
{\small
\begin{tabular}{>{\centering\arraybackslash}p{0.22\linewidth}>{\centering\arraybackslash}p{0.22\linewidth}>{\centering\arraybackslash}p{0.22\linewidth}>{\centering\arraybackslash}p{0.22\linewidth}}
\toprule
\textbf{A} & \textbf{B} & \textbf{C} & \textbf{D} \\
\midrule
1, 2 and 3 are correct & 1 and 2 only are correct & 2 and 3 only are correct & 1 only is correct \\
\bottomrule
\end{tabular}}
\par\raggedright
{\footnotesize No other combination of statements is used as a correct response.}\par
}

\begin{document}
\sethomework{AS Chemistry homework · 3.3 Metallic bonding}{19 marks}

\begin{center}
{\large\bfseries AS Chemistry homework · 3.3 Metallic bonding}\\[0.06em]
{\small (and covalent / dative / dot-and-cross)}\\[0.12em]
{\small CIE 9701 AS Chemistry · 19 marks · about 40 minutes}
\end{center}

{\small Topic: 3.3 metallic bonding; 3.4 LO1 covalent and coordinate (dative) bonding, including \ce{Al2Cl6} and \ce{NH4+}; 3.7 dot-and-cross.
Answer in English. Section A is Paper 1 style (1 mark each). Section B is written; use CIE language (electrostatic attraction, delocalised electrons, shared pair, coordinate/dative).}

\vspace{0.15em}
{\small
\begin{tabular}{@{}llr@{}}
\toprule
\textbf{Section} & \textbf{Content} & \textbf{Marks} \\
\midrule
A & Multiple choice & 8 \\
B1 & 2025 FM Paper 22 Question 2(a) & 2 \\
B2 & 2024 MJ Paper 22 Question 1(c) & 1 \\
B3 & 2024 MJ Paper 21 Question 5(a) & 1 \\
B4 & 2021 MJ Paper 22 Question 2(b) & 3 \\
B5 & 2024 MJ Paper 23 Question 4(a) & 2 \\
B6 & 2021 FM Paper 22 Question 2(e)(i) & 2 \\
\midrule
 & \textbf{Total} & \textbf{19} \\
\bottomrule
\end{tabular}}

\vspace{0.25em}
{\large\bfseries Section A --- Multiple-choice questions (8 marks)}\par

\qhead{A1.}{1}
{\small Which statement explains why calcium has a higher melting point than barium?}
\begin{enumerate}[label=\textbf{\Alph*}, leftmargin=1.7em, itemsep=0.08em]
\item Calcium cations are smaller than barium cations and have a stronger attraction to the delocalised electrons.
\item The structure of calcium is partly giant molecular.
\item There are more delocalised electrons in calcium than in barium as it has a lower ionisation energy.
\item There is greater repulsion between barium atoms as they have more complete electron shells than calcium atoms.
\end{enumerate}
\sourcefoot{2023 MJ Paper 11 Question 2}

\Needspace{16em}
\qhead{A2.}{1}
{\small Which statements help to explain the increase in melting point from sodium to aluminium?}
\begin{enumerate}[label=\arabic*., leftmargin=1.6em, itemsep=0.06em]
\item The charge on the metal ion increases.
\item There are more delocalised electrons per metal ion.
\item The radius of the metal ion decreases.
\end{enumerate}
\combokey
\sourcefoot{2021 MJ Paper 13 Question 36}

\qhead{A3.}{1}
{\small In which species does the underlined atom have an incomplete outer shell?}
\optline{\textbf{A} $\underline{\mathrm{B}}$F$_3$\qquad \textbf{B} $\underline{\mathrm{C}}$H$_3{}^{-}$\qquad \textbf{C} F$_2\underline{\mathrm{O}}$\qquad \textbf{D} H$_3\underline{\mathrm{O}}{}^{+}$}
\sourcefoot{2024 MJ Paper 12 Question 6}

\Needspace{16em}
\qhead{A4.}{1}
{\small Which species can accept a lone pair of electrons to form a coordinate (dative covalent) bond?}
\begin{enumerate}[label=\arabic*., leftmargin=1.6em, itemsep=0.06em]
\item \ce{BF3}
\item \ce{H+}
\item \ce{CH3+}
\end{enumerate}
\combokey
\sourcefoot{2021 FM Paper 12 Question 34}

\qhead{A5.}{1}
{\small When solid aluminium chloride is heated, \ce{Al2Cl6} is formed.}

{\small Which bonding is present in \ce{Al2Cl6}?}
\begin{enumerate}[label=\textbf{\Alph*}, leftmargin=1.7em, itemsep=0.08em]
\item covalent and coordinate (dative covalent)
\item covalent only
\item ionic and coordinate (dative covalent)
\item ionic only
\end{enumerate}
\sourcefoot{2024 MJ Paper 11 Question 5}

\clipq{A6.}{1}{img-006.png}
\sourcefoot{2023 ON Paper 11 Question 5}

\vspace{0.12em}
\clipq{A7.}{1}{img-007.png}
\sourcefoot{2022 FM Paper 12 Question 5}

\Needspace{16em}
\qhead{A8.}{1}
{\small Which elements form a chloride in which both covalent bonding and coordinate (dative covalent) bonding are present?}
\begin{enumerate}[label=\arabic*., leftmargin=1.6em, itemsep=0.06em]
\item Al
\item Si
\item Mg
\end{enumerate}
\combokey
\sourcefoot{2021 ON Paper 11 Question 35}

\vspace{0.35em}
{\large\bfseries Section B --- Structured questions (11 marks)}\par\vspace{0.1em}

{\normalsize\bfseries B1 · 2025 February/March Paper 22 Question 2(a) (2 marks)}\par

{\small The Group 2 elements show metallic bonding.}

\qhead{B1(a)}{2}
{\small Define metallic bonding.}
\answerlines{2}

\vspace{0.25em}
{\normalsize\bfseries B2 · 2024 May/June Paper 22 Question 1(c) (1 mark)}\par

\qhead{B2(c)}{1}
{\small Draw a labelled diagram to show the structure and bonding in sodium.}
\drawbox[4.4cm]

\vspace{0.25em}
{\normalsize\bfseries B3 · 2024 May/June Paper 21 Question 5(a) (1 mark)}\par

{\small Hydrocarbon molecules contain covalent bonds.}

\qhead{B3(a)}{1}
{\small Define covalent bond.}
\answerlines{2}

\vspace{0.25em}
{\normalsize\bfseries B4 · 2021 May/June Paper 22 Question 2(b) (3 marks)}\par

{\small Carbon monoxide, \ce{CO}, is a gas at room temperature and pressure. It contains a coordinate bond.}

\qhead{B4(b)(i)}{1}
{\small Explain what is meant by \textit{coordinate bond}.}
\answerlines{2}

\qhead{B4(b)(ii)}{2}
{\small Draw a `dot-and-cross' diagram to show the arrangement of outer electrons in \ce{CO}.}

{\small Show the electrons belonging to the C atom as $\times$.}

{\small Show the electrons belonging to the O atom as $\bullet$.}
\drawbox[4.6cm]

\vspace{0.25em}
{\normalsize\bfseries B5 · 2024 May/June Paper 23 Question 4(a) (2 marks)}\par

{\small \ce{NH3(g)} reacts with \ce{HCl(g)} to produce \ce{NH4Cl(s)}, as shown.}

\vspace{0.08em}
\centering
{\small \ce{NH3(g) + HCl(g) -> NH4Cl(s)}}
\par\raggedright

\qhead{B5(a)}{2}
{\small Draw a diagram to show the ionic, covalent and coordinate bonding present in a formula unit of \ce{NH4Cl}.}
\drawbox[4.8cm]

\vspace{0.25em}
{\normalsize\bfseries B6 · 2021 February/March Paper 22 Question 2(e)(i) (2 marks)}\par

{\small Aluminium reacts with chlorine to form aluminium chloride.}

{\small Aluminium chloride can exist as the gaseous molecule \ce{Al2Cl6(g)}. This molecule contains coordinate bonds.}

\qhead{B6(e)(i)}{2}
{\small Draw a diagram that clearly shows all the types of bond present in \ce{Al2Cl6(g)}.}
\drawbox[4.8cm]

\end{document}
